\section*{はじめに}
昨今はデータサイエンス，情報科学の領域が非常に隆盛で，コンピュータを使ってデータを分析し，経済の動向や購買行動などの予測に用いられることが広く行われている。

人の行動や考え方をどのようにデータにするかについては，当然ながら心理学には一日の長がある。
また，人が頭の中でどのような考え方のプロセスをたどるのか，それをどのように検証するのかについても，心理学はその短い歴史の中で徹底的にその技法を洗練させてきた。
このような根源的なレベルでの理論や方法論は時代が変わっても色褪せることなく，また今後ますます必要とされてくる時代になっている。

本講ではデータ解析の応用段階として，より実践的なテーマを扱う。すなわち，\textbf{心理尺度が作られる理論的背景}と，\textbf{データの背後のメカニズムを解析する方法}を知ることである。
前期配当の心理学データ解析応用1(旧カリキュラム名心理学データ解析2A)では前者の，心理尺度に関する理論的背景について学ぶ。

心理学研究法の1つとして，調査研究がある。紙とペンで回答を集めた時，回答者がある反応カテゴリにまるをつけたことが，どうして数値処理の対象になるのか。そこには数字を割り振るルールとしての「尺度化」の手続きがある。残念ながら応用的側面が発展しすぎたため，回答に数字を割り振る原理について語られることが少なくなってきてしまい，それに対する反動からか，近年改めてこの根本原理についての理解と解説が求められている。この講義では，尺度化の原理や目に見えない潜在変数を想定して分析するとはどういうことかについて，理論と演習を交えながら習得することを目指す。この理論的側面を考えるためには，どうしても線形代数・行列計算の知識が必要になってくる。線形代数については特別な事前知識は不要で，定義から改めて解説するので安心してもらいたい。

なお，2024年度から前期と後期の資料を分割して提供するようにした。通年で履修する学生が少ないことと，後期で扱うテーマとして数値シミュレーションを含むように再構成したからである。前期後期でテーマが明確になり，またスリムになった分，利用しやすくなったのではないかと思っているが，どうだろうか。

\subsection*{授業のテーマ}
データから意味のある情報を取り出すための，さまざまな分析法を習熟するにあたって，その背後にあって語られることのない「発想」の観点から理解する。
数値だけに振り回される状態から脱却し，数値を算出する数式に込められた意味について考える視点を持つ。
さらにこれらに習熟することで，どのような研究対象に対してどのような心理統計的アプローチができるかを，俯瞰的に見られるようになろう。

\subsection*{前記を通じて伝えたいポイント}
\begin{description}
	\item[尺度化とは何か] 心理学で行われるアンケート調査やその後の分析はどういう原理があって「心を測定した」といえるのか。その原理やモデルを理解して利用できるようになる。
	\item[多変量解析から何がわかるのか] 調査研究などで得られた多変量を分析することで何がわかるのか。あるいは何をしてわかったというのか。
	\item[多変量解析の基礎となる数式的原理] 多変量解析の背景にあるのは線形代数という数学であり，線形代数の基礎を学ぶことで多変量解析のメカニズムを統合的に理解できる。
\end{description}

\subsection*{その他}
授業シラバスとこの講義資料を掲載したサイト(\url{https://kosugitti.github.io/psychometrics_syllabus/})で，最新版のシラバスと授業資料，授業で用いるサンプルデータやコードの配布を行なっています。
